\chapter{Deployment}
\label{chap:deployment}

This chapter describes how to build and deploy the L3S-Offshore-2 Frontend, both locally for development and on production servers.

\section{Prerequisites}
\label{sec:prerequisites}

\subsection{Development Environment}

The following software is required for local development:

\begin{itemize}
    \item \textbf{Node.js:} Version 18.x or higher (LTS recommended)
    \item \textbf{npm:} Version 9.x or higher (comes with Node.js)
    \item \textbf{Angular CLI:} Version 17.x (\texttt{npm install -g @angular/cli})
    \item \textbf{Git:} For version control
\end{itemize}

\subsection{Production Environment}

For production deployment:

\begin{itemize}
    \item \textbf{Docker:} Version 20.x or higher
    \item \textbf{Docker Compose:} (optional, for multi-container setups)
    \item \textbf{nginx:} Included in the Docker image
\end{itemize}

\section{Local Development}
\label{sec:local-development}

\subsection{Clone and Install}

\begin{lstlisting}[language=bash, caption={Cloning and installing dependencies}]
# Clone the repository
git clone https://github.com/krausality/PNML_Frontend.git
cd PNML_Frontend

# Install dependencies
npm install
\end{lstlisting}

\subsection{Development Server}

Start the Angular development server with hot reloading:

\begin{lstlisting}[language=bash, caption={Starting the development server}]
# Start development server
ng serve

# Or with explicit host binding (for network access)
ng serve --host 0.0.0.0 --port 4200
\end{lstlisting}

The application will be available at \url{http://localhost:4200}.

\subsection{Backend Connection}

The development environment expects the Flask backend to run on port 9040. Configure the backend URL in \texttt{src/environments/environment.ts}:

\begin{lstlisting}[language=TypeScript, caption={Development environment configuration}]
export const environment = {
  production: false,
  apiUrl: 'http://localhost:9040'
};
\end{lstlisting}

\section{Production Build}
\label{sec:production-build}

\subsection{Angular Production Build}

Create an optimized production build:

\begin{lstlisting}[language=bash, caption={Creating a production build}]
# Production build with optimizations
ng build --configuration production

# Output is generated in dist/pnml-frontend/
\end{lstlisting}

The production build includes:
\begin{itemize}
    \item Ahead-of-Time (AOT) compilation
    \item Tree shaking to remove unused code
    \item Minification and uglification
    \item Content hashing for cache busting
\end{itemize}

\subsection{Build Verification}

Verify the build output:

\begin{lstlisting}[language=bash, caption={Verifying the production build}]
# Check bundle sizes
ls -la dist/pnml-frontend/browser/

# Expected output: main.js, polyfills.js, styles.css
# Total size should be under 500KB gzipped
\end{lstlisting}

\section{Docker Deployment}
\label{sec:docker-deployment}

\subsection{Dockerfile}

The application uses a multi-stage Dockerfile:

\begin{lstlisting}[language=docker, caption={Multi-stage Dockerfile}]
# Stage 1: Build
FROM node:18-alpine AS builder
WORKDIR /app
COPY package*.json ./
RUN npm ci
COPY . .
RUN npm run build -- --configuration production

# Stage 2: Serve with nginx
FROM nginx:alpine
COPY --from=builder /app/dist/pnml-frontend/browser /usr/share/nginx/html
COPY nginx.conf /etc/nginx/nginx.conf
EXPOSE 80
CMD ["nginx", "-g", "daemon off;"]
\end{lstlisting}

\subsection{Nginx Configuration}

The nginx configuration handles SPA routing:

\begin{lstlisting}[caption={nginx.conf for SPA routing}]
server {
    listen 80;
    server_name localhost;
    root /usr/share/nginx/html;
    index index.html;
    
    # SPA fallback: serve index.html for all routes
    location / {
        try_files $uri $uri/ /index.html;
    }
    
    # Cache static assets
    location ~* \.(js|css|png|jpg|jpeg|gif|ico|svg|woff|woff2)$ {
        expires 1y;
        add_header Cache-Control "public, immutable";
    }
    
    # Gzip compression
    gzip on;
    gzip_types text/plain text/css application/json application/javascript;
}
\end{lstlisting}

\subsection{Building the Docker Image}

\begin{lstlisting}[language=bash, caption={Building and running the Docker container}]
# Build the image
docker build -t l3s-offshore-frontend:latest .

# Run the container
docker run -d -p 8080:80 --name offshore-frontend l3s-offshore-frontend:latest

# Verify it's running
curl http://localhost:8080
\end{lstlisting}

\section{Server Deployment}
\label{sec:server-deployment}

\subsection{L3S Infrastructure}

The L3S research infrastructure uses a nested VPS architecture. Access requires:

\begin{enumerate}
    \item SSH to the gateway server (\texttt{gate.l3s.uni-hannover.de})
    \item SSH to the project VPS
    \item Docker commands within the VPS
\end{enumerate}

\subsection{Deployment Script}

A typical deployment workflow:

\begin{lstlisting}[language=bash, caption={Production deployment script}]
#!/bin/bash
# deploy.sh

# 1. Build and push to container registry
docker build -t registry.l3s.de/offshore/frontend:latest .
docker push registry.l3s.de/offshore/frontend:latest

# 2. On the server: pull and restart
ssh user@server << 'EOF'
  docker pull registry.l3s.de/offshore/frontend:latest
  docker stop offshore-frontend || true
  docker rm offshore-frontend || true
  docker run -d -p 80:80 --name offshore-frontend \
    registry.l3s.de/offshore/frontend:latest
EOF
\end{lstlisting}

\subsection{Docker Cleanup}

After deployment, clean up old images to save disk space:

\begin{lstlisting}[language=bash, caption={Docker cleanup commands}]
# Remove unused images (not just containers)
docker image prune -a --filter "until=24h"

# Remove stopped containers
docker container prune

# Full cleanup (use with caution)
docker system prune -a
\end{lstlisting}

\section{Environment Configuration}
\label{sec:environment-configuration}

\subsection{Production Environment}

Update \texttt{src/environments/environment.prod.ts} with production URLs:

\begin{lstlisting}[language=TypeScript, caption={Production environment configuration}]
export const environment = {
  production: true,
  apiUrl: 'https://api.offshore.l3s.de'
};
\end{lstlisting}

\subsection{Runtime Configuration}

For dynamic configuration without rebuilding, environment variables can be injected at container runtime:

\begin{lstlisting}[language=bash, caption={Runtime environment injection}]
docker run -d -p 80:80 \
  -e API_URL=https://new-api.example.com \
  l3s-offshore-frontend:latest
\end{lstlisting}

\section{Troubleshooting}
\label{sec:troubleshooting}

\subsection{Common Issues}

\begin{description}
    \item[Blank page after deployment:] Check nginx configuration and ensure \texttt{try\_files} fallback is configured for SPA routing.
    
    \item[API connection refused:] Verify the backend URL in environment configuration and check CORS settings on the backend.
    
    \item[Assets not loading:] Ensure \texttt{base href} in \texttt{index.html} matches the deployment path.
    
    \item[Container exits immediately:] Check logs with \texttt{docker logs <container-id>} and verify nginx configuration syntax.
\end{description}

\subsection{Logging}

Access container logs for debugging:

\begin{lstlisting}[language=bash, caption={Viewing container logs}]
# Follow logs in real-time
docker logs -f offshore-frontend

# Last 100 lines
docker logs --tail 100 offshore-frontend
\end{lstlisting}
